\section*{A Proof of Lemma 12}
\begin{authorclaim}{Lemma 12}
Let $T$ be a binary tree with $m$ vertices such that the $i$th vertex is labeled with $a_i\in\{0,1\}$ with the hereditary property: for each inner vertex $i$ with direct descendants $c_1$ and $c_2$, if $a_i=1$, then $a_{c_1}=1$ or $a_{c_2}=1$. We also assume that if $r$ is the root of $T$, then $a_r=1$. Assume that we have a one-sided bounded error oracle access to $a_i$ i.e. if we request a value of $a_i$ and $a_i=0$ we get 1 with probability at most $1/2$ and 0 with probability at least $1/2$; if $a_i=1$ we get 1 with probability 1. Then there exists an algorithm $A$ that with probability at least $2/3$ returns a leaf $\ell$ of $T$ with $a_\ell=1$ and makes $O(\log m)$ oracle queries to $a_1,\ldots,a_m$.
\end{authorclaim}
\begin{proof}[Proof of Lemma 12]
For a tree $F$ we denote by $|F|$ the number of nodes in $F$ and for a node $v$ of $F$ we denote by $\operatorname{Subtree}(F,v)$ the subtree of $F$ with root $v$. Let Oracle$(i)$ be the oracle function returning the correct value of $a_i$ with probability at least $9/10$. We can implement such a function using the majority vote of a constant number of initial oracle queries. Let $C$ be a constant; an appropriate value of $C$ we choose later. Consider Algorithm 2 on the following page.
We claim that at any iteration $T_i$ has the hereditary property. This is the case in the beginning and if $i$ decreases at some iteration, then the next $T_i$ was considered at an earlier iteration. Otherwise, the next $T_i$ is either a subtree of the current $T_i$ (in that case the hereditary property is clearly maintained), or is obtained by removal a subtree with 0-labeled root (here we use that the oracle has a one-sided error) from the previous $T_i$ (the hereditary property is also maintained in that case).
We first consider a variant of the algorithm that works infinitely long (i.e., $C=+\infty$) and compute the expected number of the first iteration such that $T_i$ consists of a single 1-labeled leaf of $T$. Notice that after the first such iteration the value of $T_i$ stays the same for all further iterations. We show that that the expected value is at most $C\log m$ for some constant $C$. Then by running the algorithm for $3C\lceil\log m\rceil$ iterations we obtain the required error probability by Markov's inequality.
\begin{authorclaim}{Algorithm 2: Search for 1-leaf}
\begin{algorithmic}
\State $T_0:=T$ \Comment{Initialize the tree}
\State $i:=0$
\For{$j:=1$ to $3C\lceil\log_{3/2}m\rceil$}
 \State $r:=$ root of $T_i$
 \If{Oracle$(r)=0$}
  \State $i:=\max\{0,i-1\}$ \Comment{Backtrack since the current tree may not contain a 1-leaf}
 \ElsIf{$|T_i|\ne1$}
  \State $v:=$ a centroid node of $T_i$
  \Statex \hspace{\algorithmicindent} $\triangleright$ i.e. such that $|\operatorname{Subtree}(T_i,v)|\in[\frac13|T_i|,\frac23|T_i|]$
  \If{Oracle$(v)=1$}
   \State $T_{i+1}:=\operatorname{Subtree}(T_i,v)$
  \Else
   \State $T_{i+1}:=T_i-\operatorname{Subtree}(T_i,v)$
   \Statex \hspace{\algorithmicindent} $\triangleright$ $T_{i+1}$ is obtained from $T_i$ by the deletion of $\operatorname{Subtree}(T_i,v)$
  \EndIf
  \State $i:=i+1$
 \EndIf
\EndFor
\State \Return the only node of $T_i$, if $|T_i|=1$
\end{algorithmic}
\end{authorclaim}
Let $T(j)$ denote the value of $T_i$ before the start of $j$th iteration, $i(j)$ denote $i$ at the start of $j$th iteration and $r(j)$ denote the root of $T(j)$. Notice that if $a_{r(j)}=1$, then for every $j'>j$, $T(j')$ is a subtree of $T(j)$, since the algorithm never backtracks if the true value of the roots label is 1. Hence, if $a_{r(j)}=a_{r(j')}=1$ for some $j<j'$, then $i(j)\le i(j')$.
Let us consider a sequence $j_1,j_2,j_3,\ldots$, where $j_1=0$, $j_s=\min\{j\mid a_{r(j)}=1\wedge j>j_{s-1}\wedge i(j)>i(j_{s-1})\}$, if such minimum exists.
Let us consider the iterations from $j_s$ till $j_{s+1}-1$. We consider the random variables $Y_{j_s},Y_{j_s+1},\ldots,Y_{j_{s+1}-1}$ corresponding to these iterations with the following properties:
\begin{itemize}
\item If $T(j)$ coincides with $T(j_s)$, then its root is labeled with 1. Then $Y_j=-1$ if the second oracle query returns the correct answer and $Y_j=1$ if the answer it incorrect. Notice that $\Pr[Y_j=-1]\ge9/10$.
\item If the root of $T(j)$ is labeled with zero, then $Y_j=-1$, if the first oracle query returns the correct answer (i.e. the algorithm backtracks). Otherwise, if $T(j)$ consists of a single node $Y_j=0$. Otherwise, if the root of $T(j+1)$ is labeled with 0, then $Y_j=1$. If it is labeled with 1, then $Y_j=-\infty$. Notice that $\Pr[Y_j\le-1]\ge9/10$.
\end{itemize}
Notice that, $j_{s+1}=j_s+\min\{k\mid\sum_{j=j_s}^{j_s+k-1}Y_j\le-1\}$. In order to estimate the expected value of $j_{s+1}-j_s$ we consider an auxiliary random variables $X_{j_s},X_{j_s+1},\ldots,X_{j_{s+1}-1}$, defined as
\[
X_j=\begin{cases}1,&\text{if }Y_j\ge0\\-1,&\text{if }Y_j<0.\end{cases}
\]
Notice then $\sum_{j=j_s}^{j_s+k-1}Y_j\le\sum_{j=j_s}^{j_s+k-1}X_j$. We can apply the following fact about random walks in a straight line to the random variables $X_j$:
\begin{authorclaim}{Theorem 45 (Section XII.2 of [18])}
Let $X_1,X_2,\ldots$ be a sequence of independent random variables that take value in $\{-1,1\}$. Assume that for all $i$, $\Pr[X_i=1]\le1/10$ and $\Pr[X_i=-1]\ge9/10$. Let $M$ be a random variable that equals the minimal natural number $k$ such that $\sum_{i=1}^kX_i=-1$. Then the expected value of $M$ is at most $C$, where $C\in\mathbb R$ is an absolute constant.
\end{authorclaim}
Fact 45 implies that $\mathbb E[j_{s+1}-j_s]\le C$. Then $\mathbb E[j_s]=\mathbb E[j_s-j_{s-1}+(j_{s-1}-j_{s-2})+\cdots+(j_2-j_1)+(j_1-j_0)]\le sC$. Thus, by Markov's inequality $\Pr[j_s\le3sC]\ge2/3$. Since $|T_{j_s}|\le(2/3)^s|T_{j_0}|$, the algorithm that runs for $3C\lceil\log_{3/2}m\rceil$ iterations terminates in a 1-labeled leaf with probability at least $2/3$.
\end{proof}
